\section{Proof of \Cref{prop:sufficient-cond} \label{appendix:sufficient}}

Now we turn to the proof of \Cref{prop:sufficient-cond}, restated below. Recall that $b(d, k, N; \tau)$ is the failure probability for threshold decoding at level $\tau.$

\color{blue} \begin{proposition*}
    For $\eta \in (0, 1],$ let $k \le N^\eta.$ Then for any $\epsilon > 0,$ over the regime
	$$
	d \ge (1 + \epsilon) 2(1 + \sqrt \eta)^2 k \ln N, \quad \tau \in \left [\frac 1 2 + (1 - \eta) \frac{k \ln N} d, (1 + \sqrt \eta)^{-1} \right],
	$$
	for Rademacher dictionaries it holds that
	$$
	\lim_{N \to \infty} b(d, k, N; \tau) = 0.
	$$
\end{proposition*} \color{black}

Our proof relies on the following bound for $b(d, k, N; \tau),$ valid for Rademacher dictionaries.

\begin{lemma}\label{lemma:b-bound}
For all $d, k, N \in \N$ and $\tau \in (0, 1),$
$$
b(d, k, N; \tau) \le
k \exp\left( - \frac{(1 - \tau)^2 d}{2 (k-1)} \right) + (N - k) \exp\left( - \frac{\tau^2 d} {2 k} \right).
$$
\end{lemma}

\begin{proof}
Suppose w.l.o.g. that $X = \{1, \dots, k\}$ and let $Y = FX = F_1 + \cdots + F_k.$

Define the families of events
\begin{align*}
    A_i &= [\langle F_i, Y \rangle < \tau] \quad \text{for } 1 \le i \le k, \\
    B_i &= [\langle F_i, Y \rangle \ge \tau] \quad \text{for } i > k.
\end{align*}
Let $R_t$ be a sum of $t$ independent Rademacher variables, as in \Cref{prop:rademacher-chernoff}. For $i > k,$ $\langle F_i, Y \rangle$ is distributed as $R_{kd}/d.$ Thus, by a Chernoff bound,
$$
\Pp(B_i) = \Pp(\langle F_i, Y \rangle \ge \tau) = \Pp(R_{kd} \ge d\tau) \le \exp\left(-\frac{\tau^2 d}{2k}\right).
$$
Similarly, for $i \le k,$ $\langle F_i, Y \rangle = \langle F_i, F_i \rangle + \langle F_i, \sum_{j \neq i} F_j \rangle$ is distributed like $1 + R_{(k-1)d}/d,$ and so
\begin{align*}
\Pp(A_i) &= \Pp(\langle F_i, Y \rangle < \tau) = \Pp(R_{(k-1)d} < (\tau - 1)d)\\
&\le \Pp(R_{(k-1)d} \ge (1-\tau)d) \le \exp\left(-\frac{(1-\tau)^2 d}{2(k - 1)}\right).
\end{align*}
Our conclusion follows from a union bound
$$
b(d, k, N; \tau) = \Pp\left( \bigcup_i A_i \cup \bigcup_i B_i \right) \le \sum_{i=1}^k \Pp(A_i) + \sum_{i=k+1}^N \Pp(B_i).
$$
\end{proof}

We are now prepared to complete the proof.

\color{blue}
\begin{proof}[Proof of \Cref{prop:sufficient-cond}]
Put $C = d /(k \ln N).$ Together, the assumption $k \le N^\eta$ and \Cref{lemma:b-bound} let us bound $b(d, k, N; \tau)$ by a sum of powers of $N$ with exponents depending only on $(\eta, C, \tau)$:
\begin{align*}
b(d, k, N; \tau) &\le N^\eta \exp\left(- \frac{(1 - \tau)^2 C k \ln N}{2k} \right) + N \exp\left(- \frac{\tau^2 C k \ln N}{2k} \right) \\
&= \exp\left(\left(\eta - \frac{(1-\tau)^2 C}{2}\right) \ln N\right) + \exp\left(\left(1 - \frac{\tau^2 C}{2}\right) \ln N\right) = N^{P_0} + N^{P_1}
\end{align*}
where
$$
P_0 = \eta - \frac{(1 - \tau)^2 C}{2}, \quad P_1 = 1 - \frac{\tau^2 C}{2}.
$$
For given $C$ and $\eta,$ clearly the minimum value

Since $C \ge (1 + \epsilon) 2 (1 + \sqrt \eta)^2,$
$$
asdf
$$
and so
$$
b(d, k, N; \tau) \le \exp\left(\left(\eta - \frac{\eta C}{2(1 + \sqrt{\eta})^2}\right) \ln N\right) + \exp\left(\left(1 - \frac{C}{2(1 + \sqrt{\eta})^2}\right) \ln N\right).
$$
As $N \to \infty,$ both terms vanish for any fixed $C > 2(1 + \sqrt{\eta})^2,$ which completes the proof.
\end{proof}

In the proof above, the choice of threshold $\tau = (1 + \sqrt{\eta})^{-1}$ is only justified post hoc. We now explain how to derive this expression and how it relates to the MAP threshold derived in \Cref{sec:codes}.

For a given $\eta,$ we're interested in knowing the critical value of $C$ for which some choice of threshold $\tau$ allows both terms in \eqref{two-powers} to converge to $0,$ meaning that both exponents are negative:
$$
\min_{\tau \in (0, 1)} \max \left\{ \eta - \frac{(1 - \tau)^2 C}{2}, 1 - \frac{\tau^2 C} 2 \right\} < 0.
$$
Let's write $F(\eta, \tau, C)$ for the maximum of the two exponents. The minimizer in $\tau$ is $\tau = 1/2 + (1-\eta)/C,$ at which point both exponents are equal and
$$
F(\eta, \tau, C) = \frac{1+\eta}{2} - \frac{C}{8} - \frac{(1-\eta)^2}{2C}.
$$
Solving a quadratic shows that this expression is negative so long as $C > 2(1 + \sqrt{\eta})^2.$ Substituting this critical value into our expression for the threshold $\tau$ achieving the best exponents in the upper bound gives
$$
\tau = \frac{1}{2} + \frac{1-\eta}{C} = \frac{1}{2} + \frac{1-\eta}{2(1 + \sqrt{\eta})^2} = \frac{1}{2} + \frac{(1-\sqrt{\eta})(1+\sqrt{\eta})}{2(1 + \sqrt{\eta})^2} = \frac{1}{2} + \frac{1-\sqrt{\eta}}{2(1 + \sqrt{\eta})} = \frac{1}{1 + \sqrt{\eta}}.
$$
Since $F(\eta, \tau, C)$ is decreasing in $C,$ putting $\tau = (1 + \sqrt \eta)^{-1}$ guarantees that $F(\eta, \tau, C) < 0$ for any $C$ satisfying $\min_\tau F(\eta, \tau, C) < 0.$

Finally, note that the threshold minimizing the maximum exponent $F$ of our Chernoff bounds agrees with the threshold we derived from a maximum a posteriori estimate in \Cref{sec:codes} (\Cref{eq:tau}). Indeed, substituting $C = d / (k \ln N)$ and $\eta = \ln k / \ln N$ gives
$$
\tau_{\text{Chernoff}} = \frac 1 2 + \frac{1 - \eta} C
= \frac 1 2 + \frac{k \ln N}{d} - \frac{k \ln k}{d}
= \frac 1 2 - \frac {k}{d} \ln \frac k N
$$
while, following our earlier discussion, the maximum a posteriori threshold is
\begin{align*}
\tau_{\text{MAP}} & = \frac 1 2 - \rho^{-1} \ln \frac{\Pp(X_i = 1)}{\Pp(X_i = 0)}
= \frac 1 2 - \frac{N - 1}N \frac k d \left(\ln \frac k N + \frac k N + O \left( \frac{k^2}{N^2} \right)\right) \\
& \approx \frac 1 2 - \frac k d \ln \frac k N
\end{align*}
in our regime $k \ll N,$ $k \ll d.$
